\chapter{Operadores en el dual}

Si bien ya hemos visto que las sucesiones de Sheffer se pueden definir de multitud de formas distintas, todas ellas equivalentes, aún no hemos dado ninguna forma de calcularlas ni hemos analizado la estructura algebraica que el conjunto de estas sucesiones puede poseer, asi como los casos concretos de las sucesiones de Appell y las sucesiones asociadas. En este capítulo daremos respuesta a estas dos cuestiones. Para ello, tendremos que estudiar a los operadores lineales sobre $\poly^*$. Es conveniente que pensemos en estos momentos a este como el álgebra umbral, aunque escribiremos $\poly^*$ para hacer hincapie en que los objetos que estamos manipulando son funcionales. 

Comenzamos recordando que $\poly^*$ es un espacio topológico con la topología de estabilización \ref{topologia de estabilizacion} y que esta coincicde con la topología formal, en virtud del teorema \ref{topologia formal equivalente a topologia de estabilizacion}. Esto nos permite enunciar y probar el siguiente teorema.

\begin{teorema} \label{equivalencia continuidad}
    Un operador lineal $\alpha \in \End(\poly^*)$ es continuo si y solo si verifica que para todas las sucesiones de funcionales $F_n(T)$ convergentes a 0 en la topología formal, la suma $\sum_{k=0}^\infty a_k \cdot \alpha F_k(T)$ es convergente y además
    \begin{equation*}
        \alpha \left( \sum_{k=0}^\infty a_k F_k(T) \right) = \sum_{k=0}^\infty a_k \cdot \alpha F_k(T)
    \end{equation*} 
    
\end{teorema}
\begin{demo}
    $\boxed{\impliedby}$ Sea $\{F_n(T)\}_{n \geq 0}$ una sucesión de series formales convergente a 0 en la topología formal, es decir, $\ord(F_n(T)) \rightarrow \infty$ cuando $n \rightarrow \infty$. En tal caso, por hipótesis, la suma $\sum_{k=0}^\infty a_k \alpha F_k(T)$ es convergente. Luego, por las propiedades de la topología formal, se tiene que $\ord(\alpha F_n(T)) \rightarrow \infty$ necesariamente y, por tanto, $\alpha$ es continuo (es continuo en 0 y por la linealidad es continuo en todo $\poly^*$).

    $\boxed{\implies}$ Sea $p(x) \in \poly$ y $m \geq 0$. En este caso
    \begin{equation*}
        \left\langle \alpha \sum_{k=0}^\infty a_k F_k(T) | p(x) \right\rangle = \left\langle \alpha \sum_{k=0}^m a_k F_k(T) | p(x) \right\rangle + \left\langle \alpha \sum_{k=m+1}^\infty a_k F_k(T) | p(x) \right\rangle
    \end{equation*}
    Ahora, como la sucesión formada por las $F_n(T)$ es convergente se tiene que la sucesión $\{\sum_{k=m+1}^\infty a_k F_k(T)\}_{m \geq 0}$ también es convergente. Luego, como $\alpha$ es continuo sabemos que $\{ \alpha \sum_{k=m+1}^\infty a_k F_k(T)\}_{m \geq 0}$ es convergente y, por tanto, podemos tomar $m$ lo suficentemente grande como para que el segundo término se anule y $\ord(\alpha F_k(T)) > \gr(p(x))$ para $k > m$. Así pues, obtenemos la igualdad $\langle \alpha \sum_{k=0}^\infty a_k F_k(T) | p(x) \rangle = \langle \alpha \sum_{k=0}^m a_k F_k(T) | p(x) \rangle = \langle \sum_{k=0}^m a_k \alpha F_k(T) | p(x) \rangle = \langle \sum_{k=0}^\infty a_k \alpha F_k(T) | p(x) \rangle$, que es lo que buscábamos.
\end{demo}

\begin{obs}
    Simbolizaremos como $\End_{\mathcal{C}}(\poly^*) \subset \End(\poly^*)$ al subespacio formado por los operadores lineales de $\poly^*$ que son continuos. 
\end{obs}

\section{Automorfismos y derivaciones}

Recordamos ahora la definición de automorfismo en una álgebra. El resultado de mayor relevancia para lo que nos interesa es que los automorfismos de $\poly^*$ son todos continuos. Antes de probarlo, hemos de ver un par de propiedades. 

\begin{definicion}
    Dada una $\field$-álgebra $A$ y un endomorfismo $F \in \End(A)$, se dice que $F$ es un \textbf{automorfismo} si es un morfismo de álgebras biyectivo. Simbolizaremos por $\Aut(A)$ al conjunto de automorfismos de $A$.
\end{definicion}

\begin{lema}
    Si $\alpha \in \Aut(\poly^*)$, entonces $\alpha F(T)$ es invertible (respecto al producto) si y solo si $F(T)$ es invertible (respecto al producto); es decir, $\ord(\alpha F(T)) = 0$ si y solo si $\ord(F(T)) = 0$
\end{lema}

\begin{demo}
    Si $F(T)$ es invertible con inversa $\frac{1}{F(T)}$, entonces
    \begin{equation*}
        (\alpha F(T))\left(\alpha \frac{1}{F(T)}\right)= \alpha\left(F(T) \frac{1}{F(T)}\right) = \alpha 1 = 1 \implies \frac{1}{\alpha F(T)} = \alpha \frac{1}{F(T)}
    \end{equation*}
    Recíprocamente, si $\alpha F(T)$ es invertible con inversa $\frac{1}{\alpha F(T)}$ para el producto e inversa $\alpha^{-1}$ para la composición, entonces
    \begin{equation*}
        F(T) \left(\alpha^{-1} \frac{1}{\alpha F(T)}\right) = \left(\alpha^{-1}(\alpha F(T))\right) \left(\alpha^{-1} \frac{1}{\alpha F(T)}\right) = \alpha^{-1}\left(\alpha F(T) \frac{1}{\alpha F(T)}\right) = \alpha^{-1}1 = 1
    \end{equation*}
    Con lo que $\frac{1}{F(T)} = \alpha^{-1} \frac{1}{\alpha F(T)}$ y se concluye.
\end{demo}

\begin{teorema}
    Si $\alpha \in \Aut(\poly^*)$, entonces $\ord(\alpha F(T)) = \ord(F(T))$
    para toda $F(T) \in \poly^*$
\end{teorema}

\begin{demo}
    Veamos primero que $\ord(\alpha T) = 1$. Por el lema previo, $\ord(\alpha T) \geq 1$. Si fuera mayor que 1, tendriamos que, para toda $F(T) \in \poly^*$
    \begin{equation*}
        F(T) = a_0 + T G(T) \implies \alpha F(T) = a_0 + \alpha (T \cdot G(T))
    \end{equation*}
    Ahora, $\alpha (T \cdot G(T)) = \alpha T \cdot \alpha G(T)$, por ser $\alpha$ un morfismo de álgebras, lo cual implica que $\ord(\alpha (T \cdot G(T))) \geq \ord(\alpha T) > 1$. Luego, $\ord(\alpha F(T)) \neq 1$ y no podría haber ninguna $F(T) \in \poly^*$ tal que $\alpha F(T) = T$, lo cual es imposible por la biyectividad de $\alpha$. Con lo que $\ord(\alpha T) = 1$ necesariamente. Así pues, para $k \geq 0$, $\alpha T^k = (\alpha T)^k \implies \ord(\alpha T^k) = k$

    Si ahora consideramos una $F(T) \in \poly^*$ cualquiera, con $\ord(F(T)) = k$, podemos reescibirla como $F(T) = T^k G(T)$ con $\ord(G(T)) = 0$ y por tanto llegamos a que la serie $F(T)$ cumple que $\ord(\alpha F(T)) = \ord(\alpha T^k) + \ord(\alpha G(T)) = k$
\end{demo}

\begin{teorema}\label{auto implica continuo}
    Todo automorfismo de $\poly^*$ es continuo, es decir $\Aut(\poly^*) \subset \End_{\mathcal{C}}(\poly^*)$
\end{teorema}

\begin{demo}
    Es consecuencia inmediata del teorema anterior
\end{demo}

Otra clase de operadores lineales sobre $\poly^*$ son de especial relevancia: estas son las derivaciones.

\begin{definicion}
    Dada una $\field$-álgebra $A$, se dice que un operador lineal $F \in \End(A)$ es una \textbf{derivación} si verifica que
    \begin{equation*}
        F(a \cdot b) = F(a) \cdot b + a \cdot F(b)
    \end{equation*}
    para todo $a,b \in A$. Simbolizaremos por $\Der(A)$ al conjunto de las derivaciones de $A$.
\end{definicion}

En particular, prestaremos una atención particular a las derivaciones epiyectivas, cuyo conjunto lo escribiremos como $\Der_{\mathrm{epi}}(A)$. De forma similar a lo que ocurre con los automorfismos, su propiedad más importante es que todas son continuas. Para verlo, hemos de probar primero algunas de sus propiedades. 

\begin{proposicion}\label{derivaciones nulas sobre k}
    Si $\delta \in \Der(A)$, entonces $\delta (\field) = 0$
\end{proposicion}

\begin{demo}
    $\delta 1 = \delta (1 \cdot 1) = \delta 1 + \delta 1 \implies \delta 1 = 0 \implies \delta \lambda = 0, \forall \lambda \in \field$
\end{demo}

\begin{proposicion}\label{derivaciones T^k}
    Si $\delta \in \Der(A)$ se verifica que $\delta T^k = kT^{k-1} \cdot \delta T$ para todo $k \geq 1$.
\end{proposicion}

\begin{demo}
    Para $k = 1$ la igualdad es cierta de forma evidente. Usando inducción, para $k+1$ tendremos $\delta T^{k+1} = \delta T^k \cdot T + T^k \cdot \delta T =  k T^{k} \cdot \delta T + T^k \cdot \delta T = (k+1) T^k \cdot \delta T$ 
\end{demo}

\begin{teorema}
    Si $\delta \in \Der_{\mathrm{epi}}(\poly^*)$ y $F(T) \in \poly^*$ es tal que $\ord(F(T)) \geq 1$, entonces $\ord(\delta F(T)) = \ord(F(T))-1$
\end{teorema}

\begin{demo}
    Primeramente, veamos que $\ord(\delta T) = 0$. Para ello, dada $F(T) \in \poly^*$ la escribimos como $F(T) = \sum_{k=0}^\infty a_k T^k = a_0 + TG(T)$ para cierta $G(T)$. Por tanto, por \ref{derivaciones nulas sobre k} tendremos que
    \begin{equation*}
        \delta F(T) = \delta T \cdot G(T) + T \cdot \delta G(T) \implies \ord(\delta F(T)) = \ord(\delta T \cdot G(T) + T \cdot \delta G(T))
    \end{equation*}
    Como $\delta$ es epiyectiva, podemos escoger una $F(T)$ tal que $\ord(\delta F(T)) = 0$ y $\ord(G(T)) > 0$. De tal manera que, como $\ord(T \cdot \delta G(T)) \geq 1$, necesariamente
    \begin{equation*}
        \ord(\delta T \cdot G(T)) = 0 \iff \ord(\delta T) + \ord (G(T)) = 0 \iff \ord(\delta T) = 0
    \end{equation*}
    Si ahora consideramos una serie formal cualquiera $F(T) \in \poly^*$ con $\ord(F(T)) = k$, podemos reescribirla como $F(T) = T^k G(T)$ con $\ord(G(T)) = 0$. Luego, por \ref{derivaciones T^k} llegamos a que $\ord(F(T)) = \ord(T^k \cdot \delta G(T) + kT^{k-1}G(T) \cdot \delta T)$. De aquí, usando un argumento como el de antes, obtenemos el resultado que buscábamos.
\end{demo}

De esta manera, y como vimos con los automorfismos, queda probado así la continuidad de esta clase de operadores.

\begin{teorema}\label{derivaciones epi implica continuo}
    Toda derivación epiyectiva de $\poly^*$ es continua, es decir, se verifica que $\Der_{\mathrm{epi}}(\poly^*) \subset \End_{\mathcal{C}}(\poly^*)$
\end{teorema}

\section{Operadores adjuntos}

\begin{definicion}
    Dado un operador $\Lambda \in \End(\poly)$, llamaremos \textbf{operador adjunto de $\Lambda$} al operador lineal $\Lambda^* \in \End(\poly^*)$, donde para cada $F(T) \in \poly^*$, el operador $\Lambda^*F(T)$ se define como la aplicación dada por $p(x) \mapsto \langle F(T) | \Lambda p(x) \rangle$, es decir, es aquel tal que
    \begin{equation*}
        \langle \Lambda^* F(T) | p(x) \rangle = \langle F(T) | \Lambda p(x) \rangle
    \end{equation*}
\end{definicion}

\begin{proposicion}
    \textbf{(propiedades del operador adjunto)}. Sean $\Lambda, \Omega \in \End(\poly)$ operadores cualesquiera y $\lambda \in \field$. Se verifican las siguientes propiedades.

    \begin{tabular}{@{}ll@{}}
        1. $(\Lambda + \Omega)^* = \Lambda^* + \Omega^*$ & 
        \quad 3. $(\Lambda \circ \Omega)^* = \Omega^* \circ \Lambda^*$ \\[6pt]
        2. $(\lambda \cdot \Lambda)^* = \lambda \cdot \Lambda^*$ & 
        \quad 4. $(\Lambda^{-1})^* = (\Lambda^*)^{-1}$
    \end{tabular}
\end{proposicion}

\begin{definicion}
    Dado un funcional $F(T) \in \poly^*$, se define el \textbf{operador multiplicación asociado a $F(T)$} como el operador lineal $M^*(F) \in \End(\poly^*)$ que viene definido por $G(T) \mapsto F(T) \cdot G(T)$
\end{definicion}

\begin{proposicion}
    Para toda $F(T) \in \poly^*$ se verifica que el operador asociado $M^*(F)$ es un endomorfismo continuo. Además, si $F(T) \in \F^*$, entonces $M^*(F)$ es una biyección con inversa $M^*(F)^{-1} = M(\frac{1}{F})$
\end{proposicion}

\begin{demo}
    Para la continuidad basta con aplicar el teorema \ref{equivalencia continuidad}. En el caso en el que la serie es invertible, $G(T) \xmapsto{M^*(F)} F(T) \cdot G(T) \xmapsto{M^*(\frac{1}{F})} \frac{1}{F(T)} \cdot (F(T) \cdot G(T)) = G(T) $ y viceversa prueban que es una biyección con $M^*(F)^{-1} = M(\frac{1}{F})$
\end{demo}

\begin{ej}
    \begin{enumerate}
        \item Por la proposición \ref{adjuncion serie operador y multiplicacion}, tenemos que
        \begin{equation*}
            \langle M^*(F)G(T) | p(x) \rangle = \langle G(T) | F(T)p(x) \rangle
        \end{equation*}
        Es decir, el operador adjunto de una serie formal $F(T)$ es su operador multiplicación asociado $M^*(F)$. Esto es: $F(T)^* = M^*(F)$
        \item La proposición \ref{adjuncion operador integral} se puede reescribir como
        \begin{equation*}
            \langle M^*(1/T)F(T) | p(x) \rangle = \langle F(T) | I p(x) \rangle
        \end{equation*}
        donde aquí los operadores lineales $M^*(1/T)$ e $I$ se definen como
        \begin{equation*}
            M^*(1/T)F(T) = \frac{F(T)}{T} \quad I p(x) = \int p(x) dx
        \end{equation*}
        y donde $M^*(1/T)$ solo está definido en $(\F^*)^c$. Con lo que, sobre los funcionales que no son invertibles, se verifica que $I^* = M^*(1/T)$
        \item La igualdad \eqref{adjuncion multiplicar x} implica que el operador adjunto de la multiplicación por $x$, $M$, es la derivada formal, $\partial$. Esto es: $M^* = \partial$.
    \end{enumerate}
\end{ej}

Podemos ahora construir la aplicación que mande cada operador a su operador adjunto. 
\begin{definicion}\label{aplicacion adjunta}
    Se llama \textbf{aplicación adjunta} al morfismo $\mu \colon \End(\poly) \longrightarrow \End(\poly^*)$ que viene dado por $\Lambda \mapsto \Lambda^*$
\end{definicion}

Más especifícamente, nos interesará la correspondecia existente entre un subconjunto de los operadores y los operadores continuos de $\poly^*$. En este sentido, el siguiente teorema nos muestra que este subconjunto es, de hecho, todo el espacio de operadores.

\begin{teorema}
    La aplicación adjunta restringida al espacio de los operadores lineales continuos de $\poly^*$, $\mu_{\mathcal{C}}$, es un isomorfismo con el espacio de operadores; es decir, se tiene que $\End(\poly) \backsimeq \End_{\mathcal{C}}(\poly^*)$, o lo que es lo mismo $\Im(\mu) = \End_{\mathcal{C}}(\poly^*)$ y $\mu$ es inyectiva.
\end{teorema}

\begin{demo}
    \begin{itemize}
        \item \textit{Inyectividad:} Si $\Lambda^* = 0$, entonces $\langle F(T)| \Lambda p(x) \rangle = \langle \Lambda^* F(T)| p(x) \rangle = 0$ para toda serie formal $F(T)$ y todo polinomio $p(x)$. Con lo que $\Lambda = 0$ necesariamente.
        \item \textit{Epiyectividad:} Sea $\alpha \in \End_{\mathcal{C}}(\poly^*)$. Dado un $n > 0$, por ser continuo, $\ord(\alpha T^k) > n$ para un k suficientemente grande. Así pues, podemos considerar el operador $\Lambda_{\alpha}$ que viene definido por $\Lambda_{\alpha} x^n = \sum_{k=0}^\infty \frac{\langle \alpha T^k | x^n \rangle}{k!}x^k$. Este operador cumple que $\Lambda_{\alpha}^* = \alpha$. En efecto
        \begin{align*}
            & \langle \Lambda_{\alpha}^* T^m | x^n \rangle =  \langle T^m | \Lambda_{\alpha}x^n \rangle = \sum_{k=0}^\infty \frac{\langle \alpha T^k | x^n \rangle}{k!}\langle T^m | x^k \rangle = \langle \alpha T^m | x^n \rangle\\
            & \implies \Lambda_{\alpha}^* T^m = \alpha T^m \implies \alpha = \Lambda_{\alpha}^*
        \end{align*}
        donde la última implicación se debe a la continuidad de los operadores. \qedhere
    \end{itemize}
\end{demo}

Una vez visto que la aplicación adjunta restringida $\mu_{\mathcal{C}}$ es un isomorfismo, junto a los resultados \ref{auto implica continuo} y \ref{derivaciones epi implica continuo}, nos interesaría conocer la antiimagen tanto de $\Aut(\poly^*)$ como de $\Der_{\mathrm{epi}}(\poly^*)$. Lo cierto es que estos dos conjuntos se corresponden con dos nuevas clases de operadores en $\poly^*$: estos son los \textit{operadores umbrales} y las \textit{traslaciones umbrales}, respectivamente. A su vez, estos no son más que dos subfamilias de dos clases de operadores más generales: los \textit{operadores de Sheffer} y las \textit{traslaciones de Sheffer}. Es más, estas identificaciones mediante la aplicación adjunta caracterizan a estas familias de operadores. Comenzamos estudiando y viendo este hecho para los primeros. 

%%%%%%%%%%%%%%%%%%%%%%%
%%%%%%%%%%%%%%%%%%%%%%%
%%%%%%%%%%%%%%%%%%%%%%%
%%%%%%%%%%%%%%%%%%%%%%%
%%%%%%%%%%%%%%%%%%%%%%%

%: OPERADORES DE SHEFFER
\section{Operadores de Sheffer}

\begin{definicion}
    Dada $s_n(x) \sim (G(T), F(T))$, se llama \textbf{operador de Sheffer asociado a $s_n(x)$ o a $(G(T), F(T))$} al operador $\Lambda_{G,F} \in \End(\poly)$ dado por
    \begin{equation*}
        \Lambda_{G,F} x^n = s_n(x)
    \end{equation*}
    En el caso de que $p_n(x)$ sea la sucesión asociada a $F(T)$ decimos que $\Lambda_{1, F} \equiv \Lambda_{F}$ es simplemnte el \textbf{operador umbral asociado a $p_n(x) $ o a $F(T)$}. Simbolizaremos a ambos conjuntos por $\End_{\mathrm{Sh}}(\poly)$ y $\End_{\mathrm{Umb}}(\poly)$, respectivamente.
\end{definicion}

\begin{proposicion}
    Los operadores de Sheffer son biyecciones
\end{proposicion}

\begin{demo}
        Es inmediato puesto que manda una base a otra base. 
\end{demo}

\begin{teorema}
    \textbf{(caracterización operadores umbrales)}. Se tiene la igualdad de conjuntos $\mu_{\mathcal{C}}^{-1}(\Aut(\poly^*)) = \End_{umb}(\poly)$, es decir, un operador es un operador umbral si y solo si su operador adjunto es un automorfismo de $\poly^*$. Además, para todo operador umbral $\Lambda_F \in \End_{umb}(\poly)$ y todo funcional $G(T) \in \poly^*$ se verifica que
    \begin{equation*}
        \Lambda_F^*G(T) = G(F^{-1}(T))
    \end{equation*}
\end{teorema}

\begin{demo}
    El teorema equivale a probar las inclusiones $\Aut(\poly^*) \supseteq \mu_{\mathcal{C}}(\End_{umb}(\poly))$ y $\mu_{\mathcal{C}}^{-1}(\Aut(\poly^*)) \subseteq \End_{umb}(\poly)$.

    $\boxed{\Aut(\poly^*) \supseteq \mu_{\mathcal{C}}(\End_{umb}(\poly))}$ Si la igualdad del enunciado se diera, entonces
    \begin{equation*}
        \Lambda_{F}^*(G(T)H(T)) = G(F^{-1}(T))H(F^{-1}(T)) = (\Lambda^*_{F}G(T))\cdot (\Lambda^*_{F}H(T)) 
    \end{equation*}
    y $\Lambda_{F}$ sería, por tanto, un automorfismo. Basta con ver, pues, que se tiene dicha igualdad. Por las definiciones dadas tendremos que 
\begin{align*}
    \langle \Lambda_{F}^* F^k(T)| x^n\rangle =  \langle F^k(T) | \Lambda_{F} x^n\rangle & =  \langle F^k(T) | p_n(x)\rangle = n!\delta_{n,k} = \langle T^k | x^n \rangle
\end{align*}
Luego, $\Lambda_{F}^* F^k(T) =T^k$. Y por tanto, $\Lambda_{F}^*G(F(T)) = G(T)$ para toda $G(T) \in \poly^*$, ya que los operadores adjuntos son continuos. En particular, para $G(T) = (F^{-1}(T))^k$ tenemos que $\Lambda_{F}^*T^k = (F^{-1}(T))^k$, de donde se obtiene la igualdad buscada. 

$\boxed{\mu_{\mathcal{C}}^{-1}(\Aut(\poly^*)) \subseteq \End_{umb}(\poly)}$ Sea $\alpha \in \Aut(\poly^*)$ un automorfismo. Por conservar los automorfismos el orden y ser biyecciones, habrá un único funcional delta $F(T) \in \poly^*$ tal que $\alpha F(T) = T$. En particular, por ser morfismo de algebrás será $\alpha F^k(T) = T^k$. Sea ahora $p_n(x)$ la sucesión asociada al funiconal delta $F(T)$. Tenemos que
\begin{align*}
    \langle F^k(T) | \mu_{\mathcal{C}}^{-1}(\alpha) x^n \rangle & = \langle (\mu_{\mathcal{C}}^{-1}(\alpha))^* F^k(T) | x^n \rangle = \langle \alpha F^k(T) | x^n \rangle\\
    & = \langle T^k | x^n \rangle = n! \delta_{n,k} = \langle F^k(T) | p_n(x) \rangle
\end{align*}
Por tanto, $\mu_{\mathcal{C}}^{-1}(\alpha) x^n = p_n(x)$ y, en consecuencia, $\mu_{\mathcal{C}}^{-1}(\alpha) \in \End_{\mathrm{Umb}}(\poly)$.
\end{demo}

\begin{proposicion}
    Dado $\Lambda_{G,F} \in \End_{\mathrm{Sh}}(\poly)$ se verifica que $\Lambda_{G,F} = \frac{1}{G(T)} \circ \Lambda_{F}$
\end{proposicion}

\begin{demo}
    $s_n(x) \sim (G(T), F(T)) \implies p_n(x) = G(T)s_n(x) \sim F(T)$. Así pues
    \begin{equation*}
        \Lambda_{G,F}x^n = s_n(x) = \frac{1}{G(T)} p_n(x) = \left(\frac{1}{G(T)} \circ \Lambda_{F}\right) x^n \implies \Lambda_{G,F} = \frac{1}{G(T)} \circ \Lambda_{F} \qedhere
    \end{equation*}
\end{demo}

\begin{teorema}
    \textbf{(caracterización operadors de Sheffer)}. Sea $\Lambda \in \End(\poly)$ un operador cualquiera. $\Lambda$ es un operador de Sheffer asociado a $(G(T), F(T))$ si y solo si su operador adjunto $\Lambda^*$ tiene la forma
    \begin{equation*}
        \Lambda^* = \Lambda_{F}^* \circ M^*(1/G)
    \end{equation*}
    Equivalentemente, se tiene la igualdad de conjuntos $\mu_{\mathcal{C}}(\End_{\mathrm{Sh}}(\poly)) = \Aut(\poly^*) \cup M^*(\poly^*)$ Es más, para todo operador de Sheffer $\Lambda_{G,F} \in \End_{\mathrm{Sh}}(\poly)$ y toda $H(T) \in \poly^*$ se verifica
    \begin{equation*}
        \Lambda_{G,F}^*H(T) = \frac{1}{G(F^{-1}(T))} \cdot H(F^{-1}(T)) 
    \end{equation*}
\end{teorema}

\begin{demo}
    $\boxed{\implies}$
    Se sigue directamente de la proposición anterior y las propiedades de los operadores adjuntos, de donde se ve que
    \begin{equation*}
        \Lambda_{G,F}^*H(T) = (\Lambda_{F}^* \circ M^*(1/G))H(T) = \Lambda_{F}^*\left(\frac{1}{G(T)} \cdot H(T)\right) = \frac{1}{G(F^{-1}(T))} \cdot H(F^{-1}(T)) 
    \end{equation*}

    $\boxed{\impliedby}$
    Por la proposición anterior sabemos que $\Lambda_F = G(T) \circ \Lambda_{G,F}$, lo que implica que $\Lambda_F^* = \Lambda_{G,F}^* \circ M^*(G)$. Luego, si $\Lambda^* = \Lambda_{F}^* \circ M^*(1/G)$ vemos que
    \begin{equation*}
        \Lambda^* = (\Lambda_{G,F}^* \circ M^*(G)) \circ M^*(1/G) = \Lambda_{G,F}^* \circ (M^*(G) \circ M^*(1/G)) = \Lambda_{G,F}^* 
    \end{equation*}
    de donde se llega a que $\Lambda = \Lambda_{G,F}$ usando que $\mu_{\mathcal{C}}$ es inyectiva. 
\end{demo}

De esta manera vemos que si estudiamos el comportamineto de los operadores umbrales, estudiaremos como consecuencia también a los operadores de Sheffer y viceversa. Es decir, basta con estudiar a uno de los dos. En este caso, optaremos por estudiar a los operadores umbrales en detalle, puesto que su manipulación resultará más sencilla y obtendremos como corolario una teoría para los operadores de Sheffer, aunque insistimos en que este proceso podría realizarse de forma inversa. 

Podemos ya dar una estructura algebraica tanto a los conjuntos de estos operadores como al de las sucesiones de Sheffer. 

\begin{teorema}
    $\End_{\mathrm{Umb}}(\poly)$ es un grupo bajo la ley de grupo dada por la composición de aplicaciones. Concretamente $\Lambda_F \circ \Lambda_G = \Lambda_{G \circ F} \text{ y } \Lambda_F^{-1} = \Lambda_{F^{-1}}$
\end{teorema}

\begin{demo}
    Sean $\Lambda_F, \Lambda_G \in \End_{\mathrm{Umb}}(\poly)$ dos operadores umbrales. Por su caracterización vemos que $(\Lambda_F \circ \Lambda_G)^* = \Lambda_G^* \circ \Lambda_F^* \in \Aut(\poly^*)$ ya que $\Lambda_F^*, \Lambda_G^* \in \Aut(\poly^*)$. Luego, $\Lambda_F \circ \Lambda_G \in \End_{\textrm{Umb}}(\poly^*)$ y queda así probado que es un grupo. En particular
    \begin{equation*}
        (\Lambda_F \circ \Lambda_G)^*H(T) = \Lambda_G^* H(F^{-1}(T)) = H(F^{-1}(G^{-1}(T))) = H((G \circ F)^{-1}(T)) = \Lambda_{G \circ F}^*H(T)
    \end{equation*}
    con lo que $(\Lambda_F \circ \Lambda_G)^* = \Lambda_{G \circ F}^* \implies \Lambda_F \circ \Lambda_G = \Lambda_{G \circ F}$. Análogamente se ve para $\Lambda_F^{-1}$.
\end{demo}

\begin{corolario}
    $\End_{\mathrm{Sh}}(\poly)$ es un grupo bajo la ley de grupo dada por la composición de aplicaciones. Concretamente $\Lambda_{G,F} \circ \Lambda_{H,L} = \Lambda_{G \cdot H \circ F, L \circ F} \text{ y } \Lambda_{G, F}^{-1} = \Lambda_{\frac{1}{G \circ F^{-1}}, F^{-1}}$
\end{corolario}

\begin{proposicion} Se verifica que
    \begin{enumerate}
        \item Los operadores umbrales mandan sucesiones asociadas a sucesiones asociadas
        \item El operador $\Lambda: p_n(x) \mapsto q_n(x)$, siendo $p_n(x), q_n(x)$ las sucesiones asociadas a $F(T)$ y $G(T)$, respectivamente, es un operador umbral.
    \end{enumerate}
\end{proposicion}

\begin{demo}
    \begin{enumerate}
        \item Sea $\Lambda_F \in \End_{\mathrm{Umb}}(\poly)$ y $p_n(x)$ la sucesión asociada a una serie $G(T)$. Entonces, se tiene que $\Lambda_F p_n(x) = \Lambda_F \circ \Lambda_G x^n = \Lambda_{G \circ F}x^n$, y como $\Lambda_{G \circ F}$ es un operador umbral, $\Lambda_F p_n(x)$ es una sucesión asociada.
        \item Se debe a que $\Lambda = \Lambda_F \circ \Lambda_{G^{-1}}$. \qedhere 
    \end{enumerate}
\end{demo}

\begin{corolario} Se verifica que
    \begin{enumerate}
        \item Los operadores de Sheffer mandan sucesiones de Sheffer a sucesiones de Sheffer
        \item El operador $\Lambda: s_n(x) \mapsto r_n(x)$, siendo $s_n(x)$ y $r_n(x)$ sucesiones de Sheffer, es un operador de Sheffer.
    \end{enumerate}
\end{corolario}

Aprovechando lo hecho para los operadores de Sheffer, podemos definir una operación que dota al conjunto de sucesiones de Sheffer de una estructura de grupo.

\begin{definicion}
    Sean $s_n(x) \sim (G(T), F(T))$ y $r_n(x) \sim (H(T), L(T))$ dos sucesiones de Sheffer. Se define la $\textbf{composición umbral}$ de $r_n(x) = \sum_{k=0}^n r_{n,k}x^k$ con $s_n(x)$ como el polinomio
    \begin{equation*}
        r_n(\mathbf{s(x)}) = \sum_{k=0}^n r_{n,k}s_k(x)
    \end{equation*} 
\end{definicion}

\begin{obs}
    La composición umbral se puede escribir en términos de operadores de Sheffer, pues $\Lambda_{G,F}r_n(x) = \Lambda_{G,F}[\sum_{k=0}^n r_{n,k}x^k] = \sum_{k=0}^n r_{n,k}s_k(x)$ y esto no es más que
    \begin{equation}\label{composicion umbral con operador}
        r_n(\mathbf{s(x)}) = \Lambda_{G,F}r_n(x)
    \end{equation}
\end{obs}

\begin{teorema}
    \textbf{(estructura algebraica de las sucesiones asociadas)}. $\Umb$ es un grupo \underline{no abeliano} con la ley de grupo dada por la composición umbral. En particular
    \begin{enumerate}
        \item Si $p_n(x) \sim F(T)$ y $q_n(x) \sim G(T)$, entonces $q_n(\mathbf{p(x)}) \sim G(F(T))$ 
        \item $x^n$ es la identidad para la composición umbral.
        \item Si $p_n(x) \sim F(T)$, entonces $p_n^{-1}(x) = q_n(x)$ para $q_n(x) \sim F^{-1}(T)$
    \end{enumerate}
\end{teorema}

\begin{demo}
    \begin{enumerate}
        \item Usando \eqref{composicion umbral con operador} vemos que $\Lambda_{G \circ F} x^n = q_n(\mathbf{p(x)})$ de donde se concluye.
        \item Si $p_n(x) = x^n$, tenemos que $q_n(\mathbf{p(x)}) = \sum_{k=0}^{n} q_{n,k} p_k(x) = \sum_{k=0}^{n} q_{n,k} x^k = q_n(x)$ Recíprocamente, $p_n(\mathbf{q(x)}) = \sum_{k=0}^n p_{n,k} q_k(x) = q_n(x)$
        \item Dada $p_n(x)$, por el primer apartado $x^n = \Lambda_{F^{-1} \circ F}x^n = q_n(\mathbf{p(x)})$ de donde se ve que $q_n(x) = \Lambda_{F^{-1}}x^n$ es la inversa de $p_n(x)$ y ya que esta es la sucesión asociada a $F^{-1}(T)$ se concluye. \qedhere
    \end{enumerate}
\end{demo}

\begin{corolario}
    \textbf{(estructura algebraica de las sucesiones de Sheffer)}. $\Sh$ es un grupo \underline{no abeliano} con la ley de grupo dada por la composición umbral. En particular
    \begin{enumerate}
        \item Si $s_n(x) \sim (G(T), F(T))$ y  $r_n(x) \sim (H(T), L(T))$, entonces la sucesión de la composición umbral es $r_n(\mathbf{s(x)}) \sim (G(T)H(F(T)), L(F(T)))$ 
        \item $x^n$ es la identidad 
        \item Si $s_n(x) \sim (G(T), F(T))$,entonces $s_n^{-1}(x) = q_n(x)$ para $q_n(x) \sim (G(F^{-1}(T)), F^{-1}(T))$
    \end{enumerate}
\end{corolario}

\begin{corolario}
    $\App$ es es un grupo \underline{abeliano} con la ley de grupo dada por la composición umbral.
\end{corolario}

\begin{demo}
    La conmutatividad de la composición umbral es consecuencia del apartado primero del corolario anterior, ya que si $s_n(x) \sim (G(T), T)$ y $r_n(x) \sim (H(T), T)$ esto implica que $r_n(\mathbf{s(x)}) \sim (G(T)H(T), T)$ y $s_n(\mathbf{r(x)}) \sim (H(T)G(T), T)$. Luego, obtenemos que $r_n(\mathbf{s(x)}) = s_n(\mathbf{r(x)})$.
\end{demo}

\begin{teorema}
    Para todo operador umbral $\Lambda_F \in \End_{\mathrm{Umb}}(\poly)$ y toda serie formal $G(T) \in \F$ se tiene que $\Lambda_F^* G(T) = \Lambda_{F^{-1}} \circ G(T) \circ \Lambda_F$ como operadores.
\end{teorema}

\begin{demo}
    Usando las propiedades vistas de los operadores umbrales, tenemos que
    \begin{align*}
        \langle F(T) | [\Lambda_F^* G(T)] p(x) \rangle & = \langle F(T) [\Lambda_F^*G(T)]| p(x) \rangle = \langle \Lambda_F^* [G(T) (\Lambda_F^*)^{-1}F(T)]| p(x) \rangle \\
        & = \langle G(T) (\Lambda_F^*)^{-1}F(T) | \Lambda_F p(x) \rangle 
        = \langle (\Lambda_F^{-1})^* F(T) | G(T) \circ \Lambda_F p(x) \rangle \\
        & = \langle F(T) | \lambda_F^{-1} \circ G(T) \circ \Lambda_F p(x) \rangle
    \end{align*}
    de donde vemos que $\Lambda_F^* G(T) = \Lambda_{F^{-1}} \circ G(T) \circ \Lambda_F$, como buscábamos.
\end{demo}

\begin{corolario}\label{conmutacion composicion operador umbral}
    Para todo operador umbral $\Lambda_F \in \End_{Umb}(\poly)$ y toda serie formal $G(T) \in \F$ se tiene que $\Lambda_F \circ G(F^{-1}(T)) = G(T) \circ \Lambda_F$
\end{corolario}

\begin{demo}
    Sean $F(T) \in \F$ y $p(x) \in \poly$. Tenemos que $\langle F(T) | (\Lambda_F^* G(T)) p(x) \rangle = \langle F(T) (\Lambda_F^* G(T)) |  p(x) \rangle = \langle \Lambda_F^* [G(T) (\Lambda_F^*)^{-1} F(T)] |  p(x) \rangle = \langle F(T) | \Lambda_F^{-1} G(T) \Lambda_F p(x) \rangle$, de donde se obtiene la igualdad.
\end{demo}

\begin{teorema}
    Si $p_n(x) \sim F(T)$, entonces para todo $q(x) \in \poly$ y toda $H(T) \in \F$ se tiene que $\langle H(T) | q(\mathbf{p(x)}) \rangle = \langle H(F^{-1}(T)) | q(x) \rangle $. En particular, se verifica que 
    \begin{align*}
        \langle H(T) | p_n(x) \rangle = \langle H(F^{-1}(T)) | x^n \rangle 
    \end{align*}
\end{teorema}

\begin{demo}
    Usando \eqref{composicion umbral con operador} vemos que
    \begin{equation*}
        \langle H(T) | q(\mathbf{p(x)}) \rangle = \langle H(T) |  \Lambda_{F}(q(x)) \rangle =  \langle \Lambda_{F}^*H(T) |  q(x) \rangle = \langle H(F^{-1}(T))| q(x) \rangle
    \end{equation*}
    La segunda igualdad se deduce tomando $q(x) = x^n$
\end{demo}


\begin{corolario}
    Si $s_n(x) \sim (G(T), F(T))$, entonces para todo $q(x) \in \poly$ y toda serie formal $H(T) \in \F$ se tiene que $ \langle H(T) | q(\mathbf{s(x)}) \rangle = \langle \frac{1}{G(F^{-1}(T))}H(F^{-1}(T)) | q(x) \rangle $. En particular,  se tiene que 
    \begin{align*}
        \langle H(T) | s_n(x) \rangle = \langle \frac{1}{G(F^{-1}(T))}H(F^{-1}(T)) | x^n \rangle
    \end{align*}
\end{corolario}

%%%%%%%%%%%%%%%%%%%%%%%%%%%
%%%%%%%%%%%%%%%%%%%%%%%%%%%
%%%%%%%%%%%%%%%%%%%%%%%%%%%
%%%%%%%%%%%%%%%%%%%%%%%%%%%
%%%%%%%%%%%%%%%%%%%%%%%%%%%


%: TRASLACIONES DE SHEFFER
\section{Traslaciones de Sheffer}

Una vez vistos los operadores de Sheffer, nos interesa ahora estudiar a las traslaciones de Sheffer. El desarrollo de la teoría de esta clase de operadores es muy similar a la de los operadores de Sheffer. En consecuencia, muchas de las demostraciones, por ser muy similares a las ya dadas con estos últimos, se han excluido.

\begin{definicion}
    Dada $s_n(x) \sim (G(T), F(T))$, se llama \textbf{traslación de Sheffer asociada a $s_n(x)$ o a $(G(T), F(T))$} al operador $\theta_{G,F} \in \End(\poly)$ dado por
    \begin{equation*}
        \theta_{G,F} s_n(x) = s_{n+1}(x)
    \end{equation*}
    En el caso de que $p_n(x)$ sea la sucesión asociada a $F(T)$ decimos que $\theta_{1, F} \equiv \theta_{F}$ es simplemente la \textbf{traslación umbral asociada a $p_n(x) $ o a $F(T)$}. Simbolizaremos a ambos conjuntos por $\TrasSh$ y $\TrasUmb$, respectivamente.
\end{definicion}

\begin{proposicion}
    Las traslaciones de Sheffer son biyecciones.
\end{proposicion}

\begin{teorema}
    \textbf{(caracterización de los operadores umbrales)}. Se tiene la igualdad de conjuntos $\mu_{\mathcal{C}}^{-1}(\Der_{epi}(\poly^*)) = \TrasUmb$, es decir, un operador es una traslación umbral si y solo si su operador adjunto es una derivación epiyectiva de $\poly^*$. Además, para toda traslación umbral $\theta_F \in \TrasUmb$ se verifica que $\theta_F^* = \partial_F$, es decir, es el operador de $\poly^*$ dado por
    \begin{equation*}
        \theta_F^*F^k(T) = kF^{k-1}(T)
    \end{equation*}
\end{teorema}

\begin{proposicion}
    Dada $\theta_{G,F} \in \TrasSh$ se verifica que
    \begin{equation*}
        \theta_{G,F} = M^*(1/G) \circ \theta_F \circ G(T)
    \end{equation*}
\end{proposicion}

\begin{teorema}
    \textbf{(regla de la cadena)}. Dados $\theta_F, \theta_G \in \TrasUmb$ se verifica que
    \begin{equation*}
        \theta_F = \theta_G \circ (\partial_F G(T))
    \end{equation*}
\end{teorema}

\begin{teorema}
    \textbf{(fórmulas de recurrencia)}. Sea $p_n(x) \sim F(T)$. Se verifican

    \begin{tabular}{@{}ll@{}}
        \textit{1. } $p_{n+1}(x) = M \circ \frac{1}{F'(T)}p_n(x)$ & 
        \textit{2. } $p_{n+1}(x) = M \circ \Lambda_F \circ (F^{-1}(T))' x^n$
    \end{tabular}
\end{teorema}

\begin{demo}
    \begin{enumerate}
        \item Si tomamos $G(T) = T$ se tiene que $\theta_G x^n = x^{n+1} \implies \theta_G = M$. Por el teorema anterior, vemos que $\theta_F = M \circ (\partial_F T) = M \circ \frac{1}{\partial_T F(T)} = M \circ \frac{1}{F'(T)}$. Aplicándolo a $p_n(x)$ se obtiene la fórmula primera.

        \item Se ve usando \ref{conmutacion composicion operador umbral} y el apartado anterior
        \begin{align*}
            p_{n+1}(x) = M \circ \frac{1}{F'(T)} \circ \Lambda_F x^n & = M \circ \Lambda_F \circ \frac{1}{F'(F^{-1}(T))} x^n = M \circ \Lambda_F \circ (F^{-1}(T))' x^n
        \end{align*}
    \end{enumerate}
\end{demo}

\begin{teorema}
    \textbf{(fórmula de derivación)} Dado $\theta_F \in \TrasUmb$, para todo operador $G(T) \in \End(\poly)$ se verifica que $\theta_F^* G(T) = G(T) \circ \theta_F - \theta_F \circ G(T)$. En particular, se verifica
    \begin{equation*}
        G'(T) = G(T) \circ M - M \circ G(T)
    \end{equation*}
    El lado derecho de esta igualdad se le conoce como \textbf{derivada de Pincherle}.
\end{teorema}

\begin{demo}
    Para todo $p(x) \in \poly$ y todo $G(T) \in \End(\poly)$ tenemos que
    \begin{align*}
        \langle F(T) | (\theta_F^* G(T)) p(x)\rangle & = \langle F(T)(\theta_F^* G(T)) | p(x)\rangle = \langle \theta_F^* (F(T) G(T)) - G(T) (\theta_F^* F(T)) | p(x)\rangle \\
        & = \langle F(T) | G(T) \theta_F p(x)\rangle - \langle F(T) | \theta_F G(T) p(x)\rangle \\
        & = \langle F(T) | (G(T) \theta_F - \theta_F G(T))p(x)\rangle 
    \end{align*}
    de donde se obtiene el resultado buscado.
\end{demo}

\begin{teorema}\label{recurrencia sheffer}
    \textbf{(fórmula de recurrencia para sucesiones de Sheffer)} Dada una traslación de Sheffer $\theta_{G,F} \in \TrasSh$, se verifica que $\theta_{G,F} = (M - \frac{G'(T)}{G(T)}) \circ \frac{1}{F'(T)}$
    
    En consecuencia, si $s_n(x) \sim (G(T), F(T))$ se tiene que
    \begin{equation*}
        s_{n+1}(x) = \left(M - \frac{G'(T)}{G(T)}\right) \circ \frac{1}{F'(T)}s_n(x)
    \end{equation*}
\end{teorema}

El siguiente teorema nos proporciona dos formas operacionales y directas de calcular la sucesión asociada a una serie formal $F(T)$.

\begin{teorema}\label{transferencia asociada}
    \textbf{(fórmulas de transferencia)}. Sea $p_n(x) \sim F(T)$. Se verifican las dos siguientes relaciones.
    \begin{enumerate}
        \item $p_n(x) = F'(T)(\frac{T}{F(T)})^{n+1} x^n$
        \item $p_n(x) = M \circ (\frac{T}{F(T)})^{n} x^{n-1}$
    \end{enumerate}
\end{teorema}

\begin{demo}
    Sea $G(T) = \frac{F(T)}{T}$. Veamos primero que ambas fórmulas coinciden.
    \begin{align*}
        F'(T)G^{-n-1}(T)x^n & = (TG(T))'G^{-n-1}(T)x^n = G^{-n}(T)x^n + G'(T)G^{-n-1}(T)Tx^n\\
        & = G^{-n}(T)x^n + nG'(T)G^{-n-1}(T)x^{n-1} = G^{-n}(T)x^n - (G^{-n}(T))'x^{n-1}\\
        & = G^{-n}(T)x^n - (G^{-n}(T) \circ M - M \circ G^{-n}(T)) x^{n-1} = M \circ G^{-n}(T) x^{n-1}
    \end{align*}
    donde la penúltima igualdad viene de utilizar la derivada de Pincherle sobre $G^{-n}(T)$. Basta entonces con probar que se tiene la primera de estas fórmulas. Para ello, hemos de comprobar que $q_n(x) = F'(T)G^{-n-1} (T)x^n $ cumple las condiciones de ser la sucesión asociada a $F(T)$. Usamos la caracterización \ref{sucesion asociada caracterizacion operador}.

    Sea $n \geq 1$. Usando las igualdades anteriores, vemos que
    \begin{align*}
        \langle 1| F'(T)G^{-n-1} (T)x^n \rangle = \langle 1 | G^{-n}(T) x^n - (G^{-n}(T))'x^{n-1} \rangle & = \langle 1 | G^{-n}(T) x^n \rangle - \langle 1 | \partial G^{-n}(T)x^{n-1} \rangle \\
        & = \langle G^{-n}(T) | x^n \rangle - \langle G^{-n}(T) | x^{n} \rangle = 0
    \end{align*}
    donde hemos usado que el operador adjunto de $\partial$ es $M$. Si $n = 0$, entonces $\langle 1| q_0(x) \rangle = 0$ de forma evidente. Con lo que $\langle 1| q_n(x) \rangle = \delta_{n,0} $ y la primera condición se cumple. Por otro lado
    \begin{align*}
        F(T)F'(T)G^{-n-1} (T)x^n & = TG(T)F'(T)G^{-n-1} (T)x^n = G(T)F'(T)G^{-n-1} (T) T x^n \\
        & = nF'(T)G(T)^{-n}x^{n-1} = nq_{n-1}(x)
    \end{align*} 
    y la segunda condición también se cumple. Luego, $q_n(x) = p_n(x)$ como buscábamos.
\end{demo}

Cerramos este capítulo mostrando como las fórmulas de transferencia pueden aprovecharse para obtener una fórmula con la que calcular la serie inversa, respecto a la composición, de una serie delta. Son varias las formas en las que esta fórmula puede presentarse. Damos aquí tan solo una de ellas.

\begin{teorema}
    \textbf{(fórmula de inversión de Lagrange)}. Sean $F(T) \in \F^\delta$ y $G(T) \in \F^*$. Parar todo entero $n \geq 1$ se verifica que
    \begin{equation*}
        \langle G(F^{-1}(T)) | x^n \rangle = \left\langle G'(T)\left(\frac{T}{F(T)}\right)^{n}| x^{n-1} \right\rangle
    \end{equation*} 
    En particular, para $G(T) = T$ se tiene que
    \begin{equation*}
        \langle F^{-1}(T) | x^n \rangle = \left\langle \left(\frac{T}{F(T)}\right)^{n}| x^{n-1} \right\rangle
    \end{equation*}
\end{teorema}

\begin{demo}
    Utilizando la segunda fórmula de transferencia obtenemos el resultado
    \begin{align*}
        \langle G(F^{-1}(T)) | x^n \rangle & = \langle \Lambda_F^*G(T) | x^n \rangle = \langle G(T) | \Lambda_F x^n \rangle = \langle G(T) | p_n(x) \rangle \\
        & = \left\langle G(T) | M \circ \left(\frac{T}{F(T)}\right)^{n} x^{n-1} \right\rangle = \left\langle G'(T) \left(\frac{T}{F(T)}\right)^{n}| x^{n-1} \right\rangle
    \end{align*}
    donde hemos usado nuevamente que el operador adjunto de $M$ es $\partial$.
\end{demo}

%%%%%%%%%%%%%%%%%%%%%%%%%%%
%%%%%%%%%%%%%%%%%%%%%%%%%%%
%%%%%%%%%%%%%%%%%%%%%%%%%%%
%%%%%%%%%%%%%%%%%%%%%%%%%%%
%%%%%%%%%%%%%%%%%%%%%%%%%%%